% 15_reproducibilidad.tex - Capítulo 15: Reproducibilidad y Versiones
\chapter{Reproducibilidad, Versiones y Entrega}
\label{ch:reproducibilidad}

\section{Instalación y Ejecución}
\label{sec:instalacion}

\begin{minted}{bash}
# 1. Clonar repositorio
git clone https://github.com/df1903/tactical-grid-ai
cd tactical-grid-ai

# 2. Instalar dependencias (uv provisiona Python 3.14)
uv sync

# 3. Verificar instalación
uv run pytest              # 281 tests deben pasar
uv run ruff check .        # lint
uv run mypy src            # tipos strict
\end{minted}

\section{Versiones Fijas (RE-02, RE-03, RE-04)}
\label{sec:versiones}

Todas las dependencias fijas en \code{uv.lock}:

\begin{table}[H]
\centering
\caption{Versiones fijas de dependencias}
\begin{tabular}{ll}
\toprule
\textbf{Herramienta} & \textbf{Versión} \\
\midrule
Python & 3.14 (provisionado por \code{uv}) \\
\texttt{uv} & 0.x (gestor de paquetes) \\
\texttt{pytest} & 9.1.1 \\
\texttt{ruff} & 0.16.10 (reglas E, F, I, UP, B) \\
\texttt{mypy} & 2.4.0 (--strict) \\
\texttt{minted} & (LaTeX) \\
\texttt{pgfgantt}, \texttt{pgf-pie} & (LaTeX diagrams) \\
\bottomrule
\end{tabular}
\end{table}
\label{tab:versiones}


La entrega se acompaña del repositorio Git con el código fuente, los escenarios,
los resultados y la documentación utilizados para generar estos documentos.

\section{Estados Iniciales Declarados}
\label{sec:estados-iniciales}

Todos los estados iniciales están declarados en archivos JSON de \code{scenarios/}:

\begin{itemize}
\item \code{apendice\_a.json} -- Ejemplo oficial 4$\times$4
\item \code{caso\_01\_bfs\_vs\_dfs.json} -- Caso 1
\item \code{caso\_02\_bfs\_vs\_ucs.json} -- Caso 2
\item \code{caso\_03\_obstaculo\_heuristica.json} -- Caso 3
\item \code{caso\_04\_comparacion\_heuristicas.json} -- Caso 4
\item \code{caso\_05\_beam\_k.json} -- Caso 5
\item \code{caso\_07\_decision\_adversarial.json} -- Casos 6--10 (mapa 5$\times$7)
\item \code{prueba\_15x15.json} -- Rendimiento 15$\times$15 (3 unidades/bando)
\item \code{prueba\_20x20.json} -- Mapa 20$\times$20
\item \code{escenario\_03x03} .. \code{escenario\_30x30} -- Escalabilidad
\end{itemize}

\section{Resultados Esperados vs. Obtenidos}
\label{sec:esperado-vs-obtenido}

Todos los 10 casos mínimos tienen sus valores esperados (derivados a mano) y obtenidos
(ejecución real) documentados en \texttt{cases.md} y \texttt{docs/results/}.
Los tests \code{test\_cases\_f2.py}, \code{test\_cases\_f3.py}, \code{test\_case\_6.py},
\code{test\_cases\_f4.py} asercionan \textbf{expected == obtained}.

Los archivos CSV/JSON brutos están en \texttt{docs/results/}:

\begin{itemize}
\item \code{caso\_01\_bfs\_vs\_dfs.csv/json}
\item \code{caso\_02\_bfs\_vs\_ucs.csv/json}
\item \code{caso\_03\_obstaculo\_heuristica.csv/json}
\item \code{caso\_04\_comparacion\_heuristicas.csv/json}
\item \code{caso\_05\_beam\_k.csv/json}
\item \code{casos\_07\_10\_adversarial.csv/json}
\item \code{rendimiento\_15x15.csv}, \code{rendimiento\_20x20.csv}
\item \code{apendice\_a.csv/json}
\end{itemize}

\section{Comandos de Reproducción (Evidencia RE-03, RE-04)}
\label{sec:comandos-reproduccion}

Todos los comandos exactos para regenerar los resultados están en
\texttt{user-manual.md} §6 y se resumen aquí:

\begin{minted}{bash}
R=docs/results; U="uv run python -m tactical_grid"

# Casos 1-2 (búsquedas no informadas / costos)
$U scenarios/caso_01_bfs_vs_dfs.json --algorithm BFS,DFS,UCS --verify --csv $R/caso_01_bfs_vs_dfs.csv --json $R/caso_01_bfs_vs_dfs.json
$U scenarios/caso_02_bfs_vs_ucs.json --algorithm BFS,DFS,UCS --verify --csv $R/caso_02_bfs_vs_ucs.csv --json $R/caso_02_bfs_vs_ucs.json

# Casos 3-5 (informada)
$U scenarios/caso_03_obstaculo_heuristica.json --algorithm UCS,A_STAR --heuristic H1,H2 --verify --csv $R/caso_03_obstaculo_heuristica.csv --json $R/caso_03_obstaculo_heuristica.json
$U scenarios/caso_04_comparacion_heuristicas.json --algorithm UCS,A_STAR --heuristic H1,H2 --verify --csv $R/caso_04_comparacion_heuristicas.csv --json $R/caso_04_comparacion_heuristicas.json
$U scenarios/caso_05_beam_k.json --algorithm UCS,BEAM --heuristic H1 --k 1,2,4,8 --verify --csv $R/caso_05_beam_k.csv --json $R/caso_05_beam_k.json

# Casos 6-10 (adversarial)
$U scenarios/caso_07_decision_adversarial.json --algorithm MINIMAX --depth 1,2,3 --verify --csv $R/casos_07_10_adversarial.csv --json $R/casos_07_10_adversarial.json
$U scenarios/caso_07_decision_adversarial.json --algorithm MINIMAX,ALPHA_BETA --depth 4 --verify --csv $R/casos_07_10_adversarial.csv --json $R/casos_07_10_adversarial.json
$U scenarios/caso_07_decision_adversarial.json --algorithm ALPHA_BETA --depth 3 --order default,reversed --verify --csv $R/casos_07_10_adversarial.csv --json $R/casos_07_10_adversarial.json
$U scenarios/caso_07_decision_adversarial.json --algorithm ALPHA_BETA --depth 1,2,3,4,5 --verify --csv $R/casos_07_10_adversarial.csv --json $R/casos_07_10_adversarial.json

# Rendimiento
$U scenarios/prueba_20x20.json --algorithm ALPHA_BETA --depth 2,3,4,5,6 --csv $R/rendimiento_20x20.csv
$U scenarios/prueba_15x15.json --algorithm ALPHA_BETA --depth 5,6,7 --csv $R/rendimiento_15x15.csv
$U scenarios/prueba_15x15.json --algorithm MINIMAX --depth 4,5,6 --csv $R/rendimiento_15x15_minimax.csv

# Sweep completo Apéndice A
$U scenarios/apendice_a.json --algorithm BFS,DFS,UCS,A_STAR,BEAM --heuristic H1,H2 --k 1,2,4,8 --verify --csv $R/apendice_a.csv
\end{minted}

\section{Identificación de la Entrega (RE-02)}
\label{sec:tag-entrega}

La entrega puede identificarse y reproducirse desde el repositorio del proyecto:

\begin{minted}{bash}
git describe --tags --always
\end{minted}

El comando anterior permite consultar el identificador exacto del estado entregado.

\section{Checklist Final de Entrega}
\label{sec:checklist-entrega}

\begin{longtable}{ll}
\toprule
\textbf{Ítem} & \textbf{Estado} \\
\midrule
\endfirsthead
\toprule
\textbf{Ítem} & \textbf{Estado} \\
\midrule
\endhead
\midrule
\multicolumn{2}{r}{\textit{Continúa}} \\
\endfoot
\bottomrule
\endlastfoot
Repositorio Git con aportes visibles por integrante & \checkmark \\
README.md con nombres de integrantes & \checkmark \\
Código fuente completo + instrucciones instalación/ejecución & \checkmark \\
Escenarios JSON usados + casos mínimos de prueba & \checkmark \\
Manual de usuario completo & \checkmark \\
Documento técnico completo (incluye limitaciones, reporte IA) & \checkmark \\
Escenarios JSON + 10 casos con esperado vs. obtenido & \checkmark \\
Repositorio con aportes visibles e identificador reproducible & \checkmark \\
Uso de IA declarado (herramienta, propósito, intervención, validación) & \checkmark \\
Cada integrante puede explicar cualquier parte & \checkmark \\
\bottomrule
\end{longtable}
